\section{Removing the residual of an SNC-supported divisor}
\label{sec:hironaka-divisor-unit-extension}

W{\l}odarczyk's Step~2 defines the monomial part using the irreducible
components of the boundary divisors.  Here the boundary is a finite list
of nonempty connected closed smooth hypersurfaces with simple normal
crossings.  The actual residual produced in
Theorem~\ref{hir:divisor-finite-connected-component-residual} is shown to
be the unit ideal when the original divisor is supported on that boundary.
No residual factorization or unit property is an input to this argument.
Disconnected prescribed boundary members require component refinement;
this theorem does not assign one multiplicity to such a member.

\begin{lemma}[A point on exactly one coordinate hyperplane]
\label{hir:divisor-coordinate-unit-point}
\lean{Hironaka.HironakaDivisorResidualUnit.exists_coordinate_unit_point}
\leanok
For $j\in\{1,\ldots,n\}$ and $r>0$, there is a point $z\in B(0,r)$
with $z_j=0$ and $z_i\ne0$ for every $i\ne j$.
\end{lemma}
\begin{proof}
Let $v_j=0$ and $v_i=1$ for $i\ne j$, and put
$\delta=r/(\|v\|+1)$ and $z=\delta v$.  Then $\delta>0$ and
$\|z\|=\delta\|v\|<\delta(\|v\|+1)=r$.
The asserted coordinate values follow immediately.  When $n=1$, this
construction gives $z=0$, which satisfies the same conditions.
\end{proof}

\begin{lemma}[A nonempty supported divisor has a local boundary component]
\label{hir:divisor-supported-local-component}
\uses{hir:divisor-coordinate-unit-point,trace:actual-holomorphic-coordinate-monomial-factorization,trace:actual-snc-divisor-local-monomial-generator}
\lean{Hironaka.HironakaDivisorResidualUnit.exists_member_locally_in_support}
\leanok
Let $R$ be an effective divisor on a complex $n$-manifold, and let $E$
be a finite SNC list.  Suppose $\operatorname{Supp}R\subseteq\bigcup_{W\in E}W$.
If $x\in\operatorname{Supp}R$, there is $W\in E$ containing $x$ such
that $W$ is contained in $\operatorname{Supp}R$ in a neighbourhood of $x$.
\end{lemma}
\begin{proof}
Invertibility supplies a holomorphic generator $g$ of the actual ideal
of $R$ on an open neighbourhood $U$ of $x$.  At each $y\in U$,
$y\in\operatorname{Supp}R$ if and only if $g(y)=0$: the generator itself
belongs to the ideal, and every other member is a holomorphic multiple
of it.  Choose an adapted SNC coordinate map $\psi$ with $\psi(x)=0$.
Each member of $E$ through $x$ is a single coordinate hyperplane in this
chart.  Closedness of the finitely many other members lets us shrink
the chart to avoid all of them.

The germ $G=g\circ\psi^{-1}$ has zero locus contained in the union of
coordinate hyperplanes.  The existing holomorphic coordinate
factorization theorem therefore gives
\[
 G(z)=H(z)\prod_{j=1}^n z_j^{a_j}
\]
near zero, with $H(0)\ne0$ and $a_j\in\mathbb N$.  This is a factorization
of the genuine generator, obtained by the analytic Nullstellensatz and
the primality of coordinate germs.  It is not assumed as a divisor
certificate.  Only containment of the zero locus is needed for this
local coordinate factorization.

Since $g(x)=0$, some $a_j>0$.  Otherwise every factor at zero is one,
so $g(x)=H(0)\ne0$.  This argument also rules out the support case in
ambient dimension zero, where the coordinate product is empty.
Take a sufficiently small point $z$ from
Lemma~\ref{hir:divisor-coordinate-unit-point} for this $j$.  The displayed
factorization gives $G(z)=0$, so $\psi^{-1}(z)\in\operatorname{Supp}R$.
Support containment puts that point in a member $W$ of $E$; our shrinking
ensures $x\in W$.  Its adapted coordinate equation is $z_k=0$ for a
single $k$.  All coordinates of $z$ except $j$ are nonzero, hence $k=j$.
Thus $W$ is precisely the $j$th coordinate hyperplane locally.
The same factorization vanishes at every nearby point of that
hyperplane because $a_j>0$.  Translating back to the manifold proves
the local containment of $W$ in the actual support.
\end{proof}

\begin{theorem}[No boundary component means empty support]
\label{hir:divisor-residual-empty-support}
\uses{hir:divisor-supported-local-component,hir:divisor-exponent-connected-carrier,hir:divisor-ideal-nonzero,hir:cosupport}
\lean{Hironaka.HironakaDivisorResidualUnit.support_eq_empty_of_no_component}
\leanok
Let $E$ be a finite SNC list whose members are nonempty connected closed
smooth submanifolds.  If an effective divisor $R$ satisfies
$\operatorname{Supp}R\subseteq\bigcup_{W\in E}W$ and
$W\not\subseteq\operatorname{Supp}R$ for every $W\in E$, then
$\operatorname{Supp}R=\varnothing$.
\end{theorem}
\begin{proof}
Suppose $x\in\operatorname{Supp}R$.  The preceding lemma produces a
member $W$ through $x$ locally contained in the support.  Since the
ideal of $R$ is coherent and its cosupport of mark one equals its zero
locus, this is the genuine condition
$\operatorname{IsDivNear}(I_R,W,1,x)$.
The ideal of an effective divisor has a nonzero stalk at every point.
The actual divisor exponent is locally constant on a smooth carrier,
and connectedness makes it a single integer $a$ on all of $W$.
Its divisibility characterization at $x$ gives $1\le a$.
At each $y\in W$ the same characterization now gives
$\operatorname{IsDivNear}(I_R,W,1,y)$; evaluating this neighbourhood
condition at $y$ puts $y$ in the mark-one cosupport, hence in
$\operatorname{Supp}R$.  Consequently $W\subseteq\operatorname{Supp}R$,
contrary to the hypothesis.  There can therefore be no such $x$.
\end{proof}

\begin{theorem}[The actual residual ideal is the unit ideal]
\label{hir:divisor-residual-unit}
\uses{hir:divisor-residual-empty-support}
\lean{Hironaka.HironakaDivisorResidualUnit.ideal_eq_unitIdeal_of_no_component}
\leanok
Under the hypotheses of Theorem~\ref{hir:divisor-residual-empty-support},
the complete germ ideal of $R$ equals the holomorphic unit ideal:
$I_R=\mathcal O_M$.
\end{theorem}
\begin{proof}
The support is the zero locus of the actual coherent ideal.  Empty
zero locus means that at every $x$ there is a member $b$ of the ideal
with $b(x)\ne0$.  The inverse of $b$ is holomorphic near $x$.
Every holomorphic germ $f$ is the holomorphic multiple $(f/b)b$, so it
belongs to the original ideal.  Conversely every ideal member is
holomorphic by coherence.  These two inclusions identify every stalk
with the holomorphic unit stalk, proving equality of the germ ideals.
\end{proof}

\begin{theorem}[A genuine monomial factorization of an SNC-supported divisor]
\label{hir:divisor-supported-finite-monomial}
\uses{hir:divisor-finite-connected-component-residual,hir:divisor-residual-unit,hir:divisor-finite-component-product-coherent}
\dcref{W{\l}odarczyk 2005, Step 2}
\source{wlodarczyk05-hironaka}
\lean{Hironaka.HironakaDivisorResidualUnit.exists_componentPowerProduct_eq}
\leanok
Let $E$ be a finite SNC list of nonempty connected closed smooth
hypersurfaces, and let $D$ be an effective divisor with
$\operatorname{Supp}D\subseteq\bigcup_{W\in E}W$.
There is a list $A=((W_1,a_1),\ldots,(W_s,a_s))$, with underlying list
$E$ and nonnegative integer exponents, such that
\[
 I_D=\prod_{i=1}^s I_{W_i}^{a_i}
\]
as complete germ ideals at all original points.
\end{theorem}
\begin{proof}
The existing finite connected-component extraction theorem constructs
$A$ and an actual invertible residual $R$, together with
$I_D=(\prod_i I_{W_i}^{a_i})I_R$,
$\operatorname{Supp}R\subseteq\operatorname{Supp}D$, and
$W\not\subseteq\operatorname{Supp}R$ for every $W\in E$.
The support hypothesis on $D$ therefore applies to $R$.
Theorem~\ref{hir:divisor-residual-unit} identifies $I_R$ with the unit
ideal.  Coherence of the genuine finite product and multiplication by
the unit ideal remove this factor from the full ideal equality.
The empty boundary case yields $I_D=\mathcal O_M$ by the same argument.
The exponents here are the actual exponents of the successive
intermediate divisors in the ordered extraction; no identification
with all exponents of the original divisor is asserted by this theorem.
\end{proof}
